\documentclass[11pt]{article}
\usepackage{amsmath,amssymb,hyperref}
\usepackage{xurl}
\newcommand{\OO}{\mathcal O}
\newcommand{\Nm}{\mathrm N}
\newcommand{\rad}{\operatorname{rad}}
\begin{document}

\title{The exact canonical sextic kernel for the joint detector}
\author{}
\date{}
\maketitle

All upstream references below use revision
\begin{center}
\texttt{adc7f1241b42e322a6451854ab7e4b4c146bf78a}.
\end{center}
No endpoint improvement is proved here.

\paragraph{The actual characters.}
The pinned source defines
\[
 \chi_u(K)=\eta(K)\,\mathfrak s(m^6f^4u,K),\qquad
 \mathfrak s(u,K)=\prod_{P}\chi_P(u\bmod P)^{v_P(K)}.
\]
The first identity is the ideal version of
\path{Hecke/InverseAmplificationFamily.lean:29} and
\path{Eisenstein/CompactEnergyFamilies.lean:486}.
Here \(\chi_P\) is the canonical sextic character, not an arbitrarily
chosen character of order six. In
\path{GaussSum/FiniteFourier.lean:801},
\[
 \chi_P=\mathrm{cubicChar}_P^2\,\mathrm{quadraticChar}_P.
\]
The cubic value is lifted using the reduction of the fixed Eisenstein
root \(\omega\), with
\(\overline{\mathrm{cubicChar}_P(x)}=x^{(\Nm P-1)/3}\).
The bar in this identity denotes reduction, not complex conjugation.
The resulting character has exact order six at odd primes away from
\(1-\omega\), by
\path{Arithmetic/EisensteinEmbedding.lean:471}.

The inverse detector is a M\"obius polynomial. It does not conjugate
the character. Consequently, its product with the plain polynomial
uses \(\chi_u(dn)\). A reverse fiber conjugates both factors together.
The fixed \(\eta\), \(m\), and \(f\) contribute coefficient phases and
fixed masks. They do not change the following kernel in the variable
\(u\) once those factors have been extracted.

\paragraph{The joint energy and its period.}
Let \(X=DN\) and \(s=\sigma-i\omega_0\). At a fixed common shift,
the normalized product is
\[
 I_u P_u=X^{-1/2}\sum_K c_K\chi_u(K)
       (\Nm K/X)^{-s},
\]
where
\[
 c_K=\sum_{dn=K}\mu(d)V(\Nm d/D_*)W_1(\Nm d/D)W_2(\Nm n/N).
\]
The exact finite polynomial expansion contains
\[
 X^{-1}\sum_{K,L} c_K\overline{c_L}
 (\Nm K/X)^{-s}(\Nm L/X)^{-\bar s}
 \chi_u(K)\overline{\chi_u(L)}.
\]
Each derivative of total order \(j+k\) inserts
\((-1)^ji^k\log(\Nm K/X)^{j+k}\) into the product polynomial.
These logarithms are bounded on its fixed annular support.
Different rows may have different shifts. This displayed expansion
therefore does not justify averaging rowwise witnesses at a common
shift. Uniform derivative bounds and Sobolev transport are required.

Write \(R=\rad(KL)\) and \(Q=\Nm R\).
For each \(P\mid R\), put \(q_P=\Nm P\) and
\(e_P=v_P(K)-v_P(L)\pmod6\), with representatives \(0,\ldots,5\).
Then the variable part of the row kernel is
\[
 F_{K,L}(u)=\prod_{P\mid R}
 \begin{cases}
  \chi_P(u)^{e_P},&u\notin P,\\
  0,&u\in P.
 \end{cases}
\]
The zero mask persists when \(e_P=0\).
For a local additive character \(\psi_P\), define
\[
 G_P(e,h)=\sum_{x\in\OO/P}(\chi_P^e)(x)\psi_P(hx).
\]
The zeroth power here is a multiplicative character, whose value at
zero is zero. It is not the scalar convention \(0^0=1\).
Under CRT, for an additive character expressed in local coordinates,
\[
 \widehat F_{K,L}(h)=\prod_{P\mid R}G_P(e_P,h_P).
\]
The exact local classification is
\[
 G_P(e,0)=
 \begin{cases}q_P-1,&e=0,\\0,&e\ne0,\end{cases}
 \qquad
 G_P(0,h)=-1\quad(h\ne0).
\]
For \(e\ne0\), \(h\ne0\), and primitive \(\psi_P\),
\[
 G_P(e,h)=(\chi_P^e)(h)^{-1}\tau_P(e),
 \qquad |\tau_P(e)|^2=q_P.
\]
These local identities are already formalized upstream in
\path{Arithmetic/EisensteinEmbedding.lean:400--458}.
An additive character with a minus sign replaces \(h\) by \(-h\).
One must also include the actual CRT frequency scaling for the
chosen global trace character.

\paragraph{Exceptional pairs are sixth equivalent.}
The complete zero frequency is exactly
\[
 \sum_{u\bmod R}F_{K,L}(u)=
 \begin{cases}
  \prod_{P\mid R}(q_P-1),&v_P(K)\equiv v_P(L)\pmod6\text{ for all }P,\\
  0,&\text{otherwise}.
 \end{cases}
\]
The exceptional class consists of
\(K=A B^6\), \(L=A C^6\), where \(A\) has prime exponents in
\(\{0,1,\ldots,5\}\). Distinct ideals can be exceptional.
For example, \(K=P Q^7\) and \(L=P^7Q\) have the same complete
kernel on units modulo \(PQ\), although \(K\ne L\) when \(P\ne Q\).
Thus a claim of complete cancellation for every distinct pair is false.
This does not falsify the proposed joint moment saving. The total
resonant energy still requires summation with normalized coefficients.

The new file \path{lean/JointKernelArithmetic.lean} proves the
actual ideal identity
\[
 \mathfrak s(u,K A^6)=\mathfrak s(u,K)
          \mathbf 1_{(A,(u))=1}.
\]
It also proves equality of the two sixth-twisted factors on their
common unit set and canonical active and masked zero-mode identities.
These statements preserve the support masks explicitly.

\paragraph{What completion does not establish.}
If \(K,L\) are coprime squarefree ideals with norm comparable to
\(X\), then \(Q\asymp X^2\), all their local exponents are active,
and every primitive nonzero complete Fourier kernel has norm
\(\sqrt Q\asymp X\). A smooth row cutoff of area \(U\) has roughly
\(Q/U\) significant dual frequencies. Bounding the completed sum
term by term gives order \(\sqrt Q\), up to cutoff constants.
Near \(r+m=1.123\), this is \(U^{1.123}\), which exceeds the
trivial per-pair bound \(U\).
Therefore complete zero-mode cancellation and local Gauss bounds alone
do not yield the required joint energy
\[
 \sum_u |\partial_\sigma^j\partial_{\omega_0}^k(I_uP_u)|^2
 \ll \mathrm{height}^J U^{1+(5/6)(r+m-1)-\eta+\rho},
 \qquad j+k\le2,\quad\eta>0.
\]
The new estimate must control the weighted sum over \(K,L\) and
nonzero dual frequencies. A detector fiber indicator also prevents
direct use of the complete mean. One can dominate its nonnegative
energy by the full family before derivative transport, but that step
does not supply the missing cancellation.

There is an exact finite formulation of the missing estimate.
For a finitely supported row weight \(w\), set
\[
 w_R(x)=\sum_{u\equiv x\bmod R}w(u),\qquad
 \widehat w_R(h)=\sum_{x\bmod R}w_R(x)\psi_R(-hx).
\]
Finite Fourier inversion gives
\[
 \sum_u w(u)F_{K,L}(u)
   =Q^{-1}\sum_{h\bmod R}\widehat w_R(h)
                    \prod_{P\mid R}G_P(e_P,h_P).
\]
Write
\(a_K=c_K\eta(K)\mathfrak s(m^6f^4,K)(\Nm K/X)^{-s}\).
Once the resonant pairs are bounded, the remaining target is a bound
for the real part of
\[
 \frac1X\sum_{K\not\sim_6 L}a_K\overline{a_L}
 \frac1{\Nm\rad(KL)}
 \sum_{h\ne0}\widehat w_{\rad(KL)}(h)
                \prod_{P\mid\rad(KL)}G_P(e_P,h_P)
\]
with exponent \(1+\frac56(r+m-1)-\eta+\rho\).
The period and the Fourier transform of the row weight depend on the
pair. This dependence must remain in a proposed bilinear estimate.

Even at one prime, restricting a complete sum can remove its
cancellation. For the subset \(S=\{u:\chi_P(u)=1\}\),
\(\sum_{u\in S}\chi_P(u)=(q_P-1)/6\), whereas the complete sum is
zero. This example diagnoses an invalid inference for arbitrary
subsets. It makes no claim that the detector fibers equal \(S\).

\paragraph{Exact computation.}
Run
\begin{verbatim}
uv run research/joint_kernel.py \
  --checkpoint research/joint_kernel.checkpoint.json
\end{verbatim}
The script models both primes above split rational primes and inert
residue fields \(\mathbf F_p[\omega]/(\omega^2+\omega+1)\).
It computes the canonical cubic and quadratic values by finite-field
powering. It stores Fourier sums in
\(\mathbf Z[\zeta_6,\zeta_p]\), imposing both cyclotomic relations.
Gauss norm identities and complete joint exponent sums are checked
with integers. No complex floating-point comparisons are used.
The run checked 1,716 local Fourier kernels, including their canonical
Gauss phases, and 1,440 direct joint exponent products in ten residue
fields. It also records the size of the subset where the canonical
character equals one.
Checkpoints retain the script hash, pinned revision, field labels,
and a digest of every exact Fourier kernel.
Resume with \texttt{--resume}. A changed script requires a fresh run.

\paragraph{Four shifts with the actual exponent labels.}
For \(\theta=\chi_P^e\), \(1\le e\le5\), define
\[
 C_\theta(a,b,c,d)=\sum_{x\in\OO/P}
 \theta(x+a)\theta(x+b)\overline{\theta(x+c)\theta(x+d)}.
\]
Let \(m_t=\#\{a,b\}\cap\{t\}-\#\{c,d\}\cap\{t\}\), where
the counts include multiplicity. The exceptional divisor condition is
\(e m_t\equiv0\pmod6\) for every distinct shift \(t\).
On this exceptional variety the complete sum is exactly
\(q_P-\#\{a,b,c,d\}\).
For \(e=1,2,4,5\), the exceptional condition says that the numerator
and denominator multisets agree. For \(e=3\), it additionally includes
\(a=b\), \(c=d\), even if \(a\ne c\).
Excluding only the usual paired shifts is therefore insufficient for
a uniform square-root bound on the actual local exponent family.

Three collision identities are exact. If \(b\ne d\) and
\(\theta\ne1\),
\[
 C_\theta(0,b,0,d)=-1-\theta(b)\overline{\theta(d)},
 \qquad |C_\theta(0,b,0,d)|\le2.
\]
For \(a\ne c\),
\[
 C_\theta(a,a,c,c)=
 \begin{cases}q_P-2,&\theta^2=1,\\-1,&\theta^2\ne1.\end{cases}
\]
These statements are proved in \texttt{JointKernelArithmetic.lean}
using the exact shifted-character correlation from the pinned source.

For \(t\ne0,1\), the substitution \(y=x/(x+1)\), followed by
\(z=(t-1)y/t\), gives the three-root Jacobi identity
\[
 C_\theta(0,0,1,t)=
   \theta\!\left(\frac{t}{(t-1)^2}\right)
       J(\theta^2,\theta^{-1})-1.
\]
The omitted point \(y=1\) has value one. The pole \(x=-1\) has
value zero, so it contributes no additional term.
If \(\theta^2\ne1\), both Jacobi characters and their product are
nonprincipal, and the exact Jacobi product theorem gives
\(|J(\theta^2,\theta^{-1})|^2=q_P\).
Hence the three-root bound is \(\sqrt{q_P}+1\).
The new Lean file proves this general Jacobi norm identity from
\path{Mathlib.NumberTheory.JacobiSum.Basic}.
The displayed M\"obius substitution is proved here algebraically and
checked by the exact script, but is not formalized in Lean.

For four distinct roots, the proposed bound
\[
 |C_\theta(a,b,c,d)|\le2\sqrt{q_P}+1
\]
is an unproved input in this work. No general finite-field Weil bound
for this rational character sum was found in the imported Lean
theorems. The script tests every translation and scaling class over
the ten actual fields, for all five nonzero exponent labels.
Its checkpoint labels this uniform bound as observed, not proved.
The computations give two exact counterexamples to deleting the
additive constant in this bound. Above the split rational prime
\(19\), either embedding of \(\omega\) gives
\[
 C_{\chi_P}(0,1,2,3)=-9,\qquad 9^2>4\cdot19.
\]
In the inert field \(\mathbf F_{121}\),
\[
 C_{\chi_P^3}(0,1,\omega,3+5\omega)=-23,
 \qquad 23=2\sqrt{121}+1.
\]
All four shifts are distinct in each example.
The second example attains the proposed bound exactly.
The script checked 93,210 normalized four-shift cases in total.

\paragraph{The Type II reorganization and its cost.}
At a fixed shift write
\[
 \alpha_d=\mu(d)V(\Nm d/D_*)W_1(\Nm d/D)(\Nm d/D)^{-s},
 \qquad \beta_n=W_2(\Nm n/N)(\Nm n/N)^{-s}.
\]
Fixed base-character phases and masks can be included in these
coefficients. The full Type II energy before grouping is
\[
 \frac1{DN}\sum_{d_1,d_2,n_1,n_2}
 \alpha_{d_1}\overline{\alpha_{d_2}}
 \beta_{n_1}\overline{\beta_{n_2}}
 \mathcal K_w(d_1n_1,d_2n_2).
\]
Both M\"obius factors remain present.
For the nonzero Fourier remainder, replace this full row kernel by
\[
 \mathcal K_w^{\ne0}(K,L)=
   \frac1{\Nm\rad(KL)}\sum_{h\ne0}
      \widehat w_{\rad(KL)}(h)\prod_{P\mid\rad(KL)}G_P(e_P,h_P).
\]
For nonresonant pairs this agrees with the full variable row kernel,
since the complete zero mode vanishes. This statement keeps the
fixed base phases in the coefficients \(\alpha_d\), \(\beta_n\).
They do not guarantee sign cancellation in every sector.
If \(d_1,d_2\) are distinct prime ideals, their product of M\"obius
values is \(+1\). Taking \(n_1,n_2\) to be two further distinct
prime ideals gives the full period \(\Nm\rad(d_1d_2n_1n_2)\asymp
(DN)^2\). Thus the largest periods occur in a sector where these
two M\"obius signs supply no cancellation.
In the full Type II energy, the diagonal strips \(d_1=d_2\) and
\(n_1=n_2\) reduce to the
marginal plain and inverse energies with bounded annular weights.
They are already order \(U^{1+\rho}\) when those marginal estimates
hold. Removing these strips does not change the cardinality exponent
of the remaining pairs.

One exact Cauchy step for the full positive joint energy is instructive.
For nonnegative \(w\), put
\[
 \begin{aligned}
 A_u&=\sum_d\alpha_d\chi_u(d),\qquad
 B_u=\sum_n\beta_n\chi_u(n),\\
 T(d_1,d_2)&=\sum_u w(u)|B_u|^2
                \chi_u(d_1)\overline{\chi_u(d_2)}.
 \end{aligned}
\]
Then
\[
 E^2\le\frac{(\sum_d|\alpha_d|^2)^2}{D^2N^2}
          \sum_{d_1,d_2}|T(d_1,d_2)|^2,
\]
and the latter sum is exactly
\[
 \sum_{u,v}w(u)w(v)|B_u|^2|B_v|^2
 \left|\sum_d\chi_u(d)\overline{\chi_v(d)}\right|^2.
\]
This step replaces the signed M\"obius coefficients by their squared
norm. It has not created a M\"obius saving. Reversing the roles of
\(d\) and \(n\) retains the M\"obius factors in \(|A_u|^2|A_v|^2\),
but requires a new bound for the correlated plain sum in the last
factor.

There is also an explicit cost even under a stronger hypothetical
input. Suppose every nonexceptional incomplete pair kernel were
bounded by \(U^{1/2+\rho}\). Summing its absolute value over
\(D^2N^2\) pairs, then dividing by \(DN\), gives
\[
 E_{\mathrm{off}}\ll DN\,U^{1/2+\rho}
       =U^{r+m+1/2+\rho}.
\]
At \((r,m)=(0.71499156,0.40839014)\), this exponent is
\(1.62338170\). The required exponent before a saving is
\(1+\frac56(r+m-1)=1.10281808\).
The missing gain from the weighted average is therefore at least
\[
 (r+m+1/2)-\left(1+\frac56(r+m-1)\right)
       =\frac{r+m}{6}+\frac13=0.52056362.
\]
Cauchy with only a mean square pair-kernel bound of order \(U\)
has the same cost. Thus even hypothetical square-root cancellation
for each pair does not prove the needed mixed moment by itself.

The four-shift sums above are the complete local kernels produced by
additive differencing of a fixed residue character. They are not a
replacement for the actual Type II expression with varying periods
\(\rad(d_1n_1d_2n_2)\).
Applying them globally requires a proved reorganization that preserves
those periods, their CRT scaling, and the signed coefficient average.
That bridge and the required weighted saving remain unproved.

\paragraph{A complete differencing kernel on the full-conductor sector.}
Fix pairwise coprime squarefree \(d_1,d_2,n_1,n_2\), and put
\(Q_0=d_1d_2n_1n_2\). The pair character is now a primitive product
\(\Theta_{Q_0}=\prod_{P\mid Q_0}\chi_P^{e_P}\), where each
\(e_P\) equals \(1\) or \(5\).
For fixed additive translates \(a,b,c,d\), its complete four-shift
correlation factors exactly as
\[
 \mathcal C_{Q_0}(a,b,c,d)=
  \prod_{P\mid Q_0}C_{\chi_P^{e_P}}(a_P,b_P,c_P,d_P).
\]
The quadratic double-root exception is absent in this sector.
Let \(E\) be the squarefree product of those primes \(P\mid Q_0\)
for which either
\((a_P,b_P)=(c_P,d_P)\) or
\((a_P,b_P)=(d_P,c_P)\).
The exact exceptional values, together with the unproved generic
local bound, give the conditional estimate
\[
 |\mathcal C_{Q_0}(a,b,c,d)|
    \le3^{\#\{P:P\mid Q_0\}}
           \sqrt{\Nm Q_0\,\Nm E}.
\]
Keeping \(E\) is necessary when translates coincide modulo only
some of the conductor primes.

For two shifts, no new Weil input is needed. CRT and the exact
nonprincipal shifted-character identity give
\[
 \sum_{x\bmod Q_0}\Theta_{Q_0}(x+h)
                       \overline{\Theta_{Q_0}(x)}=
 \prod_{P\mid Q_0}
 \begin{cases}q_P-1,&h\in P,\\-1,&h\notin P.\end{cases}
\]
At shifts coprime to \(Q_0\), this complete correlation has norm
one. Applying it to an incomplete row region still requires a
controlled completion or differencing argument. It does not remove
the dependence of \(Q_0\) on the original quadruple.

\paragraph{Grouping by the reduced sextic label while retaining masks.}
Write each nonzero ideal uniquely as \(K=A_K B_K^6\), with the
prime exponents of \(A_K\) in \(\{0,\ldots,5\}\).
For the following finite identity assume that each sixth part
\(B_K\) is supported, as required by the upstream sixth-mask theorem.
In particular this holds in the prime-product sector below, where
the sixth part is one. The exact identity is
\[
 \sum_K a_K\mathfrak s(u,K)
 =\sum_A\mathfrak s(u,A)
       \sum_{K:A_K=A}a_K\mathbf1_{(B_K,(u))=1}.
\]
This finite grouping identity is now proved in
\path{group_finite_ideal_rows_by_reduced_label_with_sixth_mask}.
The inner coefficient depends on \(u\). Dropping that dependence
would discard the sixth-power zero mask.
An equivalent grouping has two labels \((A,T)\), where
\[
 T=\frac{\rad(B_K)}{\gcd(\rad(B_K),\rad(A))}.
\]
For that grouping put
\[
 C_{A,T}=\sum_{K:(A_K,T_K)=(A,T)}a_K.
\]
Then
\[
 \sum_K a_K\mathfrak s(u,K)
 =\sum_{A,T}C_{A,T}\mathfrak s(u,A)
                    \mathbf1_{(T,(u))=1}.
\]
Removing primes already dividing \(A\) from \(T\) is valid because
\(\mathfrak s(u,A)\) itself vanishes there.

For two labels \((A,T)\), \((A',T')\), the reduced multiplicative
conductor is
\[
 Q_6=\prod_{v_P(A)-v_P(A')\not\equiv0\bmod6}P.
\]
The union zero mask is still supported on
\(R=\rad(AA'TT')\). Its passive part is \(Z=R/Q_6\).
Thus the pair kernel is a primitive character modulo \(Q_6\),
multiplied by the coprimality mask at \(Z\), and its full period is
\(R=Q_6Z\). Grouping by \(Q_6\) alone does not determine the kernel.
Pairs with \(Q_6=1\) remain the sixth-equivalent exceptional pairs.

\paragraph{Singleton prime products.}
The actual source profile satisfies
\(\operatorname{supp}W\subset(1/2,2)\), by
\path{DetectorDictionaryInverseUniform.positiveAnnular_support_sharp}
in \path{Dictionary/InverseUniformSource.lean:14}.
Hence every product coefficient has \(\Nm K<4X\).
If \(\Nm A\ge X\) and \(K=A B^6\) lies in this support, then
\(B=1\). Indeed, a nonunit nonzero integral ideal has norm at least
two, so \(\Nm K\ge64\Nm A\ge64X\).
This implication is proved in
\path{exclude_nontrivial_sixth_part_from_full_norm_annulus}.
Every squarefree product \(PQ\) with \(\Nm P\ge D\) and
\(\Nm Q\ge N\) therefore occupies a singleton reduced-label class.
Its additional mask label is \(T=1\).

Here are the precise source-window hypotheses for a square-mass
lower bound. Choose fixed closed intervals
\(J_I=[a_I,b_I]\), \(J_P=[a_P,b_P]\subset[1,2)\), of positive length,
on which the real source profiles obey
\(W_1\ge w_I>0\), \(W_2\ge w_P>0\).
Assume \(D>4N\), \(N>2\), and that
\(V(\Nm P/D_*)=1\) on the chosen inverse-prime interval.
At the stated crossing the last condition holds for all sufficiently
large \(U\) whenever \(r\le t_* -\kappa\), with \(\kappa>0\), because
the actual cutoff is one on \([0,1]\).
These positivity hypotheses hold for the actual
\texttt{positiveAnnular}: for \(1\le y<2\) it equals
\texttt{sieveBump}(y), since the cutoff at \(2y\) is zero.
The bump has radii one and two and is positive inside its outer ball.
Continuity gives a positive minimum on every closed subinterval of
\([1,2)\). This assertion uses the concrete bump, not only a support
condition for an arbitrary profile.

For primes \(P,Q\) in these respective windows, the only divisors of
\(PQ\) are \(1,P,Q,PQ\). Separation of the annuli makes the inverse
weight vanish at \(1,Q,PQ\). Consequently
\[
 c_{PQ}=-W_1(\Nm P/D)W_2(\Nm Q/N),\qquad
 |c_{PQ}|^2\ge w_I^2w_P^2.
\]
The factorization alternatives and this coefficient evaluation are
proved in \path{classify_prime_product_factor_pair} and
\path{evaluate_separated_prime_product_coefficient}.
The separation also makes the map \((P,Q)\mapsto PQ\) injective.
No M\"obius convolution cancellation occurs in this sector.

The prime-counting input is available for these actual Eisenstein
ideals in the pinned upstream. In
\path{Hecke/RayPrimes.lean:38},
\path{HeckeRayPrimes.classCoeff_ratio_tendsto} proves the
ideal Mangoldt mass limit in a fixed principal ray class, with
positive density \(A=(\operatorname{rayCard}M)^{-1}\).
For a smooth nonnegative real bump \(\varphi\) supported inside
either positive window, with positive integral, the theorem
\path{PNT.AnnularPrimeMass.primeAnnularSum_log_div_tendsto}
in \path{PrimeCounting/AnnularPrimeMass.lean:335} gives
\[
 \frac{\log x}{x}\sum_{P\in\mathcal C}
       \varphi(\Nm P/x)\longrightarrow A\int\varphi>0.
\]
Its hypotheses include differentiability of \(\varphi\), continuity
of its derivative, zero endpoint values and a positive lower endpoint.
The displayed conclusion implies at least \(c x/\log x\) primes in
the window, since \(\varphi\) is bounded. A fixed finite set of bad
primes can be deleted once the window moves beyond their norms.
These upstream results supply the counting input. The asymptotic
lower bound below is deduced here and is not a new Lean theorem.

It follows that, for the pure source coefficients and all sufficiently
large \(D,N\), with fixed windows, ray class and deleted prime set,
\[
 \sum_K|c_K|^2\ \ge\sum_{P,Q}|c_{PQ}|^2
       \gg \frac{DN}{\log D\log N}.
\]
The same lower bound holds after grouping by \((A,T)\), because
these classes are singletons. Bounded common shifts multiply each
singleton coefficient by a factor bounded above and below on the
chosen windows. Fixed character phases have modulus one after the
fixed bad primes are deleted. Thus the obstruction also applies to
the unmarked shifted coefficient square norm.
It rules out a uniform bound \(X^{1-\varepsilon}\), for any fixed
\(\varepsilon>0\), obtained solely by cancellation in these grouped
coefficients. It does not obstruct the desired joint energy: the
normalized diagonal contribution remains only order \(U\), up to
logarithms, below the target exponent.

\paragraph{What remains after reduced-conductor grouping.}
The nonresonant grouped energy still asks for a bound on
\[
 \frac1X\operatorname{Re}\sum_{(A,T)\not\sim_6(A',T')}
 C_{A,T}\overline{C_{A',T'}}
 \sum_u w(u)\mathfrak s(u,A)\overline{\mathfrak s(u,A')}
                 \mathbf1_{(TT',(u))=1}.
\]
Here nonresonance means \(A\ne A'\), since both labels are sixth
reduced. The weight, derivative transport and fixed coefficient
phases must still be included as before. The singleton prime-product
sector contains coprime pairs with \(\Nm Q_6\asymp X^2\), \(Z=1\),
and positive products of the two M\"obius signs. Reduced-conductor
grouping therefore neither removes the largest conductors nor
forces coefficient sign cancellation on them.

At \((r,m)=(0.71499156,0.40839014)\), the required exponent is
\[
 1.10281808-\eta+\rho,\qquad \eta>0.
\]
The absolute use of even a hypothetical \(\sqrt U\) pair bound
still has exponent \(1.62338170\), with the previously computed
gap \(0.52056362\). The counting lower bound prevents recovering
that gap from a coefficient square-norm power saving alone.
A bound for the weighted average of the pair characters, with their
active conductors and passive masks, remains the missing input.

The exact script also checks this grouping for indicator annuli.
In each run the square masses before and after masked grouping agree,
and the maximum group size is one.
\begin{center}
\begin{tabular}{rrr}
\((D,N)\) & Square mass & Prime products with \(c_{PQ}=-1\)\\
\hline
\((32,8)\) & 14 & 12\\
\((128,16)\) & 229 & 110\\
\((256,32)\) & 770 & 246\\
\((512,64)\) & 2885 & 1140\\
\((1024,64)\) & 5988 & 2085
\end{tabular}
\end{center}
These are arithmetic checks using the actual split and inert prime
ideals away from two and three. They use indicator profiles and do
not simulate the smooth detector energy or the sixth-free row family.

\end{document}
